\section{Experimental details}
\label{appendix:experimental-details}

\subsection{Tabular density estimation}
\label{appendix:experimental-details-tabular-density-estimation}

Model selection is performed using the standard validation splits for these datasets.
We clip the norm of gradients to the range $ [-5, 5] $, and find this helps stabilize training.
We modify MAF by replacing permutations with invertible linear layers.
Hyperparameter settings are shown for coupling flows in \cref{tab:experimental-details-density-estimation-coupling} and autoregressive flows in \cref{tab:experimental-details-density-estimation-autoregressive}.
We include the dimensionality and number of training data points in each table for reference. 
For higher dimensional datasets such as Hepmass and BSDS300, we found increasing the number of coupling layers beneficial. 
This was not necessary for Miniboone, where overfitting was an issue due to the low number of data points.

\begin{table*}[!ht]
	\caption{Hyperparameters for density-estimation results using coupling layers in \cref{sec:tabular-density-estimation}.}
	\label{tab:experimental-details-density-estimation-coupling}
    \vspace*{-0.1in}
	\begin{center}
		\begin{small}
			\begin{sc}
				\begin{tabular}{lccccc}
					\toprule
					& Power & Gas & Hepmass & Miniboone & BSDS300 \\ 
					\midrule
					Dimension & 6 & 8 & 21 & 43 & 63 \\
					Train data points & 1,615,917 & 852,174 & 315,123 & 29,556 & 1,000,000 \\
					\midrule
					Batch size & 512 & 512 & 256 & 128 & 512 \\
					Training steps & 400,000 & 400,000 & 400,000 & 200,000 & 400,000 \\
					Learning rate & 0.0005 & 0.0005 & 0.0005 & 0.0003 & 0.0005 \\
					Flow steps & 10 & 10 & 20 & 10 & 20 \\
					Residual blocks & 2 & 2 & 1 & 1 & 1 \\
					Hidden features & 256 & 256 & 128 & 32 & 128 \\
					Bins & 8 & 8 & 8 & 4 & 8 \\
					Dropout & 0.0 & 0.1 & 0.2 & 0.2 & 0.2 \\
					\bottomrule
				\end{tabular}
			\end{sc}
		\end{small}
	\end{center}
\end{table*}

\begin{table*}[!h]
	\caption{Hyperparameters for density-estimation results using autoregressive layers in \cref{sec:tabular-density-estimation}.}
	\label{tab:experimental-details-density-estimation-autoregressive}
	\vspace*{-0.1in}
	\begin{center}
		\begin{small}
			\begin{sc}
				\begin{tabular}{lccccc}
					\toprule
					& Power & Gas & Hepmass & Miniboone & BSDS300 \\ 
					\midrule
					Dimension & 6 & 8 & 21 & 43 & 63 \\
					Train data points & 1,615,917 & 852,174 & 315,123 & 29,556 & 1,000,000 \\
					\midrule
					Batch size & 512 & 512 & 512 & 64 & 512 \\
					Training steps & 400,000 & 400,000 & 400,000 & 250,000 & 400,000 \\
					Learning rate & 0.0005 & 0.0005 & 0.0005 & 0.0003 & 0.0005 \\
					Flow steps & 10 & 10 & 10 & 10 & 10 \\
					Residual blocks & 2 & 2 & 2 & 1 & 2 \\
					Hidden features & 256 & 256 & 256 & 64 & 512 \\
					Bins & 8 & 8 & 8 & 4 & 8 \\
					Dropout & 0.0 & 0.1 & 0.2 & 0.2 & 0.2 \\
					\bottomrule
				\end{tabular}
			\end{sc}
		\end{small}
	\end{center}
\end{table*}

\begin{table*}[!h]
	\caption{Validation log likelihood (in nats) for UCI datasets and BSDS300, with error bars corresponding to two standard deviations.}
	\label{tab:uci-results-validation}
        \vspace*{-0.1in}
	\begin{center}
		\begin{small}
			\begin{sc}
				\begin{tabular}{lccccc}
					\toprule
					Model & POWER & GAS & HEPMASS & MINIBOONE & BSDS300 \\ 
					\midrule
					RQ-NSF (C) & $ 0.65 \pm 0.01 $ & $ 13.08 \pm 0.02 $ & $ -14.75 \pm 0.06 $ & $ \hphantom{0}{-9.03} \pm 0.43 $ & $ 172.51 \pm 0.60 $  \\
					RQ-NSF (AR) & $ 0.67 \pm 0.01 $ & $ 13.08 \pm 0.02 $ & $ -13.82 \pm 0.05 $ & $ \hphantom{0}{-8.63} \pm 0.41 $ & $ 172.5 \pm 0.59 $ \\
					\bottomrule
				\end{tabular}
			\end{sc}
		\end{small}
	\end{center}
	\vspace{-1.0em}
\end{table*}

\subsection{Improving the variational autoencoder}
We use the Adam optimizer \citep{kingma:2014:adam} with default hyperparameters, annealing an initial learning rate of $ 0.0005 $ to $ 0 $ using a cosine schedule \citep{loshchilov:2016:sgdr} over $ 150{,}000 $ training steps with batch size $ 256 $. 
We use a `warm-up' phase for the KL divergence term of the loss, where the multiplier for this term is initialized to $ 0.5 $ and linearly increased to $ 1 $ over the first $ 10\% $ of training. 
This modification initially reduces the penalty incurred by the approximate posterior in deviating from the prior, and similar schemes have been shown to improve VAE training dynamics \citep{rezende:2018:taming}. 
Model selection is performed using a held-out validation set of $ 10{,}000 $ samples for MNIST, and $ 20{,}000 $ samples for EMNIST\@.

We use 32 latent features, and residual nets use 2 blocks, with 64 latent features for coupling layers, and 128 latent features for autoregressive layers. 
Both coupling and autoregressive flows use 10 steps. 
As with the tabular density-estimation experiments, we modify IAF \citep{Kingma:2016:iaf} and MAF \citep{Papamakarios:2017:maf} by replacing permutations with invertible linear layers using an LU-decomposition. 
All NSF models use 8 bins. 
The encoder and decoder architectures are set up exactly as described by \citet{nash2019aem}, and are similar to those used in IAF \citep{Kingma:2016:iaf} and NAF \citep{Huang:2018:naf}. 

Conditioning the approximate posterior distribution $ q(\bfz \g \bfx) $ follows a multi-stage procedure. 
First, the encoder computes a context vector $ \bfh $ of dimension $ 64 $ as a function of the input $ \bfx $. 
This vector is then mapped to the mean and diagonal covariance of a Gaussian distribution in the latent space. 
Then, $ \bfh $ is also given as input to the residual nets in each of the flow's coupling or autoregressive layers, where it is concatenated with the input $ \bfz $, mapped to the required number of hidden features, and also used to modulate the additive update of each residual block with a sigmoid gate. 
We found this scheme to work well across experiments.

\subsection{Generative modeling of images}

For image-modeling experiments we use a Glow-like model architecture introduced by \citet[Section~3]{Kingma:2018:glow}.
This involves stacking multiple steps for each level in the multi-scale architecture of \citet{Dinh:2017:rnvp}, where each step consists of an actnorm layer, an invertible $1\times1$ convolution and a coupling transform. For our RQ-NSF (C) model,
we make the following modifications to the original Glow model: we replace affine coupling transforms with rational-quadratic coupling transforms, we go back to residual convolutional networks as used in RealNVP \cite{Dinh:2017:rnvp}, and we use an additional $1\times 1$ convolution at the end of each level of transforms. The basline model is the same as RQ-NSF (C), except that it uses affine coupling transforms instead of rational-quadratic ones.
For CIFAR-10 experiments we do not factor out dimensions at the end of each level, but still use the squeezing operation to trade spatial resolution for depth.

For all experiments we use 3 residual blocks and batch normalization \citep{ioffe:2015:batchnorm} in the residual networks which parameterize the coupling transforms. We use 7 steps per level for all experiments, resulting in a total of 21 coupling transforms for CIFAR-10, and 28 coupling transform for ImageNet64 (Glow models used by \citet{Kingma:2018:glow} use 96 and 192 affine coupling transforms for CIFAR-10 and ImageNet64 respectively).

We use the Adam \citep{kingma:2014:adam} optimizer with default $\beta_1$ and $\beta_2$ values. An initial learning rate of 0.0005 is annealed to 0 following a cosine schedule \citep{loshchilov:2016:sgdr}. We train for $100{,}000$ steps for 5-bit experiments, and for $200{,}000$ steps for 8-bit experiments.
To track the performance of our models, we split off 1\% of the training data to use as a development set. Due to the resource requirements of the experiments, we perform a limited manual hyper-parameter exploration. Final values are reported in \cref{tab:experimental-details-images}.

We use a single NVIDIA Tesla P100 GPU card per CIFAR-10 experiment, and two such cards per ImageNet64 experiment. Training for $200{,}000$ steps takes about $5$ days with this setup.

\begin{table*}[h!]
	\caption{Hyperparameters for generative image-modeling results in \cref{sec:gen_image_modeling}.}
	\label{tab:experimental-details-images}
	\vspace*{-0.1in}
	\begin{center}
		\begin{small}
			\begin{sc}
				\begin{tabular}{lcccccc}
					\toprule
					Dataset && Batch Size & Levels & Hidden Channels & Bins & Dropout \\ 
					\midrule
					CIFAR-10 & 5-bit & 512 & 3 & 64 & 2 & 0.2\\
					& 8-bit & 512 & 3 & 96 & 4 & 0.2\\
					ImageNet64 & 5-bit & 256 & 4 & 96 & 8 & 0.1\\
					& 8-bit & 256 & 4 & 96 & 8 & 0.0\\
					\bottomrule
				\end{tabular}
			\end{sc}
		\end{small}
	\end{center}
\end{table*}
